%% ============================================================================
%% APPENDIX M — Supplementary Technical and Data Screening Tables
%% Tables moved from main chapters to reduce body length while preserving
%% full methodological transparency.
%% ============================================================================

%% When included as a standalone appendix, uncomment the \chapter line below:
%% \chapter{Supplementary Technical and Data Screening Tables}


\noindent\textbf{Purpose.} This appendix is provided to support transparency, traceability, reproducibility, and examination defensibility. The material contained herein supplements, but does not replace, the primary findings reported in Chapters~1--5. No major conclusions rely exclusively on appendix content.

% [SPACE-FIX] clearpage removed to eliminate empty page
\section{Supplementary Technical and Data Screening Tables}
\label{app:supplementary_tables}

Data screening tables, model architecture comparisons, and EEG frequency band reference are provided in the Supplementary Volume.

%% MOVED TO SUPPLEMENTARY VOLUME: Data Screening and Model Tables
\iffalse

\subsection{Data Screening Tables}

The following four tables document the data quality assurance procedures
applied prior to hypothesis testing. They are referenced from
Section~\ref{sec:thematic_analysis} of Chapter~4.

\noindent Table~\ref{tab:missing_data} summarises missing data assessment and treatment by dataset for appendix review.

{\tablestripes\tablefontsize
\needspace{20\baselineskip}
\begin{xltabular}{\textwidth}{@{} L c L L @{}}
\caption{Missing Data Assessment and Treatment by Dataset}\label{tab:missing_data} \\
\toprule
\thead{Dataset} & \thead{Missing (\%)} & \thead{Cause} & \thead{Treatment} \\
\midrule
\endfirsthead
\endhead
\bottomrule
\endlastfoot
KAU ASD       & 0.0 & --- (real-time QA protocols) & None required \\
PhysioNet PD  & 2.3 & Electrode displacement (T5, T6) & Per-channel median imputation \\
DEAP Stress   & 0.0 & --- (controlled lab environment) & None required \\
SAM40 Stress  & 1.1 & Bluetooth connectivity gaps & Median imputation \\
BCI-IV 2a     & 0.0 & --- (competition-grade curation) & None required \\
TCIA Cancer   & 3.7 & Missing slice metadata / DICOM errors & Records excluded \\
\end{xltabular}}
\vspace{0.3em}\noindent\textit{Note.} All rates below the 5\% threshold for multiple imputation. TCIA records excluded (not imputed) because volumetric feature extraction requires complete data.

\noindent Table~\ref{tab:app_outlier_detection_supp} summarises outlier detection for appendix review.

{\tablestripes\tablefontsize
\needspace{20\baselineskip}
\begin{xltabular}{\textwidth}{@{} L c L c @{}}
\caption{Outlier Detection: IQR-Flagged Rates and Treatment by Dataset}\label{tab:app_outlier_detection_supp} \\
\toprule
\thead{Dataset} & \thead{IQR-Flagged (\%)} & \thead{Treatment} & \thead{Final Removal (\%)} \\
\midrule
\endfirsthead
\endhead
\bottomrule
\endlastfoot
KAU ASD       & 1.4 & Artefact epochs removed; physiological retained & 1.4 \\
PhysioNet PD  & 0.9 & Winsorised & 0.9 \\
DEAP Stress   & 1.8 & Muscle/blink artefacts removed & 1.8 \\
SAM40 Stress  & 1.6 & Same as DEAP & 1.6 \\
BCI-IV 2a     & 2.1 & Electrode-pop epochs removed & 2.1 \\
TCIA Cancer   & 0.3 & Winsorised & 0.3 \\
\end{xltabular}}
\vspace{0.3em}\noindent\textit{Note.} All removal rates below the 5\% threshold. Physiologically plausible outliers retained with winsorisation; only artefact-contaminated epochs removed.

\noindent Table~\ref{tab:normality_testing} summarises normality testing for appendix review.

{\tablestripes\tablefontsize
\needspace{20\baselineskip}
\begin{xltabular}{\textwidth}{@{} l r r l L @{}}
\caption{Normality Testing: Shapiro--Wilk Results per Dataset}\label{tab:normality_testing} \\
\toprule
\thead{Dataset} & \thead{W Statistic} & \thead{p-value} & \thead{Normal?} & \thead{Action Taken} \\
\midrule
\endfirsthead
\endhead
\bottomrule
\endlastfoot
KAU ASD & 0.934 & 0.021 & No & Non-parametric tests (Wilcoxon) used \\
PhysioNet PD & 0.967 & 0.183 & Yes & Parametric tests permissible \\
DEAP Stress & 0.912 & 0.004 & No & Bootstrap CI for inference \\
SAM40 Stress & 0.945 & 0.048 & No & Non-parametric + bootstrap \\
BCI-IV 2a & 0.891 & 0.012 & No & Friedman test for 4-class comparison \\
TCIA Cancer & 0.978 & 0.312 & Yes & Parametric tests permissible \\
\end{xltabular}}
\vspace{0.5em}\noindent\textit{Note.} Normality assessed at $\alpha$ = .05. Datasets failing the Shapiro--Wilk test were analysed using non-parametric alternatives or bootstrap-based inference. The fold-level accuracy distributions (reported in the main hypothesis testing sections) exhibited approximate normality for the ASD domain, permitting paired \textit{t}-tests with Bonferroni correction as the primary inferential method.

\noindent Table~\ref{tab:class_balance} summarises class balance assessment and mitigation for appendix review.

{\tablestripes\tablefontsize
\needspace{20\baselineskip}
\begin{xltabular}{\textwidth}{@{} L L c L @{}}
\caption{Class Balance Assessment and Mitigation}\label{tab:class_balance} \\
\toprule
\thead{Domain} & \thead{Distribution} & \thead{Ratio} & \thead{Mitigation} \\
\midrule
\endfirsthead
\endhead
\bottomrule
\endlastfoot
ASD   & 47\% ASD / 53\% TD & 1:1.13 & Stratified split \\
PD    & 50\% PD / 50\% HC & 1:1.00 & Balanced by design \\
Stress & 40--43\% high / 57--60\% low & 1:1.3--1.5 & Stratified split + $F_1$ reporting \\
BCI   & 25\% per class (4-class) & 1:1:1:1 & Balanced by design \\
Cancer & 46\% malignant / 54\% benign & 1:1.17 & Stratified split \\
\end{xltabular}}
\vspace{0.3em}\noindent\textit{Note.} CV = cross-validation. Subject-level splitting prevents epoch-level data leakage.

\clearpage

\addcontentsline{toc}{section}{Model Specification Tables (Chapter~3)} % [TOC-appendix-section-insertion]
\section*{Model Specification Tables (Chapter~3)}

The following two tables document the model architectures and hyperparameter
configurations used across the experimental pipeline. They are referenced from
Section~\ref{sec:analysis_techniques} of Chapter~3.

\noindent Table~\ref{tab:model_architecture_comparison} summarises model architecture comparison for appendix review.

{\tablestripes\tablefontsize
\needspace{20\baselineskip}
\begin{xltabular}{\textwidth}{@{} L L L L L @{}}
\caption[Model Architecture Comparison: Parameters, Inputs,]{Model Architecture Comparison: Parameters, Inputs, Advantages, and Domain Deployment}\label{tab:model_architecture_comparison} \\
\toprule
\thead{Model} & \thead{Parameters} & \thead{Input} & \thead{Advantage} & \thead{Domains Applied} \\
\midrule
\endfirsthead
\endhead
\bottomrule
\endlastfoot
\multicolumn{5}{l}{Classical Baselines} \\
\midrule
SVM & N/A (kernel) & Feature vector & High-dim margin optimisation & All ASD \\
\addlinespace
Random Forest & $\sim$10K trees & Feature vector & Built-in importance; nonlinear & All ASD \\
\addlinespace
XGBoost & $\sim$5K boosted & Feature vector & Regularised boosting; tabular data & All ASD \\
\addlinespace
LDA & $\sim$100 & CSP features & Low complexity; BCI standard & BCI, PD \\
\midrule
\multicolumn{5}{l}{Deep Learning} \\
\midrule
EEGNet & $\sim$2.5K & ($C$, $T$) & Lightweight; interpretable; depthwise conv & ASD, PD, BCI \\
\addlinespace
DeepConvNet & $\sim$300K & ($C$, $T$) & EEG-optimised deep extraction & BCI, Stress \\
\addlinespace
1D-CNN & $\sim$150K & ($T$,) & Temporal convolutions on raw signal & ASD, PD, Stress \\
\addlinespace
2D-CNN & $\sim$500K & ($H$, $W$, 3) & Spatial features from scalograms/images & All EEG + Cancer \\
\addlinespace
CNN-LSTM & $\sim$1M & ($C$, $T$) & Spatial + temporal + attention & ASD, PD, Stress \\
\addlinespace
3D-CNN & $\sim$25M & ($D$, $H$, $W$) & Volumetric medical imaging & Cancer \\
\end{xltabular}}
\vspace{0.5em}\noindent\textit{Note.} $C$ = channels; $T$ = time points; $H$ = height; $W$ = width; $D$ = depth (volumetric). SVM = support vector machine; LDA = linear discriminant analysis; CSP = common spatial pattern; CNN = convolutional neural network; LSTM = long short-term memory. Parameter counts are approximate and vary with input dimensions. EEGNet's compact architecture ($\sim$2.5K parameters) is specifically designed for small EEG datasets (BCI: $n = 9$), while CNN-LSTM's larger capacity ($\sim$1M parameters) is applied to datasets where temporal dynamics are critical (ASD, PD, Stress).

\noindent Table~\ref{tab:hyperparameter_config} summarises model hyperparameter configuration for appendix review.

{\tablestripes\tablefontsize
\needspace{20\baselineskip}
\begin{xltabular}{\textwidth}{@{} L L L L L L @{}}
\caption{Model Hyperparameter Configuration}\label{tab:hyperparameter_config} \\
\toprule
\thead{Parameter} & \thead{SVM} & \thead{RF} & \thead{XGBoost} & \thead{CNN-LSTM} & \thead{EEGNet} \\
\midrule
\endfirsthead
\endhead
\bottomrule
\endlastfoot
Kernel / Estimators & RBF & 500 trees & 300 rounds & --- & --- \\
\addlinespace
$C$ / Max depth & 1.0 & 20 & 6 & --- & --- \\
\addlinespace
$\gamma$ / Learning rate & auto & --- & 0.1 & 0.001 & 0.001 \\
\addlinespace
Batch size & --- & --- & --- & 32 & 32 \\
\addlinespace
Epochs & --- & --- & --- & 100 & 100 \\
\addlinespace
Early stopping & --- & --- & 10 rounds & 15 epochs & 15 epochs \\
\addlinespace
Optimiser & --- & --- & --- & Adam & Adam \\
\addlinespace
Dropout & --- & --- & --- & 0.5 & 0.25 \\
\addlinespace
Weight decay & --- & --- & L2 = 0.01 & L2 = $10^{-4}$ & L2 = $10^{-4}$ \\
\end{xltabular}}
\vspace{0.5em}\noindent\textit{Note.} SVM = support vector machine; RF = random forest; RBF = radial basis function; $C$ = regularisation parameter; $\gamma$ = kernel bandwidth; L2 = L2 regularisation (weight decay). Dashes indicate parameters not applicable to the given model family. All deep learning models use Adam optimiser with $\beta_1 = 0.9$, $\beta_2 = 0.999$. Hyperparameters were selected via grid search within inner cross-validation folds.

\clearpage

\addcontentsline{toc}{section}{EEG Reference Table (Chapter~1)} % [TOC-appendix-section-insertion]
\section*{EEG Reference Table (Chapter~1)}

The following table documents the five canonical EEG frequency bands
and their clinical associations relevant to this dissertation.
It is referenced from the Signal Processing Theories section of Chapter~1.

\noindent Table~\ref{tab:eeg_frequency_bands} summarises eeg frequency bands and clinical significance for appendix review.

{\tablestripes\tablefontsize
\needspace{20\baselineskip}
\begin{xltabular}{\textwidth}{@{} p{2.0cm} p{2.0cm} p{2.0cm} p{8.5cm} @{}}
\caption{EEG Frequency Bands and Clinical Significance}\label{tab:eeg_frequency_bands} \\
\toprule
\thead{Band} & \thead{Range (Hz)} & \thead{Amplitude ($\mu$V)} & \thead{Clinical Association} \\
\midrule
\endfirsthead
\endhead
\bottomrule
\endlastfoot
Delta ($\delta$) & 0.5--4 & 20--200 & Deep sleep stages (N3), diffuse brain injury, metabolic encephalopathy; elevated delta power is a non-specific marker of cortical dysfunction \\
\addlinespace
Theta ($\theta$) & 4--8 & 5--100 & Drowsiness, meditative states, working memory encoding; elevated frontal theta is a documented ASD biomarker~\cite{bosl2018eeg}; theta/beta ratio is used in ADHD assessment \\
\addlinespace
Alpha ($\alpha$) & 8--13 & 10--50 & Relaxed wakefulness with eyes closed; posterior alpha suppression correlates with PD tremor severity~\cite{geraedts2018eeg}; alpha asymmetry is a stress indicator \\
\addlinespace
Beta ($\beta$) & 13--30 & 5--30 & Active cognitive processing, motor planning, sustained attention; excessive beta synchronisation in basal ganglia--cortical loops is a PD hallmark; beta desynchronisation drives BCI motor imagery classification \\
\addlinespace
Gamma ($\gamma$) & 30--100 & 2--10 & High-frequency cognitive binding, perceptual integration, cross-modal processing; reduced gamma coherence is associated with ASD social cognition deficits; gamma-band features support BCI control signal decoding \\
\end{xltabular}}
\vspace{0.5em}\noindent\textit{Note.} Amplitude ranges are approximate and vary with electrode placement, subject age, and recording conditions. Hz = Hertz; $\mu$V = microvolts. The RGAIG preprocessing pipeline extracts power spectral density (PSD) features from all five bands, enabling the multi-scale CNN architecture to learn domain-specific discriminative patterns.

\fi %% END MOVED TO SUPPLEMENTARY VOLUME

\clearpage
